\section{Stage the complete result before publishing it}
Six successful block calls are needed for one complete macroblock result.
Writing each result straight into the caller's buffer exposes a partially updated
macroblock if a later call fails. Instead, reserve private space for all six
records, populate those records, and copy them to the caller after success.
The dependency's output pointer must refer to the private record while it runs.

A reference frame plan uses 144 bytes for a request copy, 64 for one reusable
block request and 1152 for staged results, plus saved state. Record the ownership
of each region. A nested predictor can allocate its own frame below the current
stack pointer; that space must not overlap your live records. The public harness
provides 16\,384 bytes of private stack and checks its outer guards. Its size
permits other valid frame arrangements.

\begin{figure}[htbp]
\centering
\begin{tikzpicture}[every node/.style={draw,rounded corners,align=center,font=\small},node distance=6mm,alt={The macroblock routine first validates its request and all six planes. It then forms six prediction and residual records in private storage. If a block predictor fails, it returns failure while retaining the caller's old output. After all six succeed, it copies the complete records into the caller's output and returns success.}]
\node (admit) {validate request and all six planes};
\node[below=of admit] (stage) {produce six records\\in private storage};
\node[below=of stage] (publish) {after all succeed, copy 1152 bytes\\and return success};
\draw[-{Latex}] (admit)--(stage);
\draw[-{Latex}] (stage)--(publish);
\end{tikzpicture}
\caption{Private staging keeps a failed macroblock from publishing partial records.}
\end{figure}

The guarantee is evaluated when the function returns. A 1152-byte copy consists
of multiple stores; it does not provide hardware-atomic visibility to concurrent
readers. With the caller's valid mapped output and ordinary single-threaded
simulator use, all successful records are available at the return boundary.

\section{Observe one complete raw macroblock}
Use \path{construction/h263-macroblock-construction/main.asm}. Carry in your completed
eight-by-eight predictor and retain one ordinary main entry. The supplied
macroblock entry and predictor initially return development status four.
Prepare a new file from your current working source:
\begin{verbatim}
python3 construction/h263-macroblock-construction/prepare.py \
    --source my-macroblock.asm --mode 1 --dx -2 --dy -2 \
    --output worked-macroblock.asm
\end{verbatim}
Preparation installs two full raw QCIF pictures and selects the mode and signed
displacement. It leaves every instruction byte unchanged and refuses an existing
output. The ordinary caller initializes all six plane pointers at runtime, using
25344-byte Y and 6336-byte chroma plane sizes. Its current origin is $(16,32)$.
The supplied program entry first rounds its stack pointer down to a sixteen-byte
boundary. MARS's legacy initial stack value is not a promise of that alignment.
The entry uses a full-width mask of $-16$ before the first nested call; its
terminal syscall exits the program after the observation. Keep an already aligned
caller's stack unchanged across the callable routine itself.

Unlike the earlier block helper, it does not calculate individual block pointers
or fractional phases for the target routine.

For the generated pictures, the component sums are:
\begin{center}
\tagpdfsetup{table/header-rows={1}}
\begin{tabular}{rrrrr}
\textbf{Displacement} & \textbf{Four Y blocks} & \textbf{Cb} & \textbf{Cr} & \textbf{Total}\\
$(-2,-2)$ & 2976 & 0 & 0 & 2976\\
$(-1,-1)$ & 5208 & 232 & 348 & 5788\\
$(0,0)$ & 7872 & 480 & 720 & 9072
\end{tabular}
\end{center}
The first row agrees with the selected luma SAD, but that agreement alone does
not prove correct chroma handling: its two chroma residuals happen to be zero.
The odd displacement exercises half-sample chroma phases and distinguishes the
two chroma planes. Its Cb and Cr sums differ because their sample values differ.

The completed default worked caller prints \texttt{2976} and exits zero.
At \texttt{macro\_caller\_return}, before its first instruction, all six records
are already published and return registers are zero and 2976. The caller has not
yet copied those registers into its named status and total fields. Try the odd
and zero displacements as separate prepared files; their totals are 5788 and 9072.
This observation still uses raw pictures to isolate integration. The eventual
encoder must supply its locally reconstructed reference planes.
